\section{Related Work in Music Programming}\label{sec:related-work}

Existing music programming systems differ in live versus offline targeting, in whether composition is separated from
synthesis, and in how temporal structure is expressed.
Several of the systems below belong to the live-coding tradition, in which a performer writes and edits program text
on stage as an established performance practice in its own right~\cite{Collins_McLean_Rohrhuber_Ward_2003}.
This section surveys representative languages in turn; for each, it describes how the system approaches composition,
the niche it shares with Motivo, and the gap that remains relative to this thesis.
The research-gap section below then synthesizes those gaps into the position Motivo occupies.

\topic{Sonic Pi.} Embeds Ruby for education and live performance~\cite{Aaron_Blackwell_2013,Aaron_2016}, wrapping
SuperCollider behind \texttt{play} and \texttt{sleep}.
Its temporal model is imperative: a thread-local clock advances as statements run, so onsets are relative to execution
time and long-range form is expressed through live program order rather than a pre-materialized timeline.
Like Motivo, it exposes composition as readable, textual source rather than hand-placed events, and it shares the
conviction that program structure is a natural way to describe music.
Aaron emphasizes classrooms and on-stage improvisation, where immediate feedback outweighs freezing a finished
score~\cite{Aaron_2016}, and every \texttt{play} assumes a running synthesizer; no compile pass validates a complete
song and emits a portable artifact.
For Motivo, the gap is performance coupling: offline, deterministic symbolic compilation is not a first-class outcome,
and the rules that generated a passage are recoverable only by re-reading and re-running the live program.

\topic{TidalCycles and Strudel.} Represent music as time-varying patterns composed
functionally~\cite{McLean_Wiggins_2010}.
Time is cyclic rather than imperative: patterns describe events within repeating cycles, combined and transformed
while a performance clock runs, which suits reshaping rhythm and harmony mid-set.
Their pattern algebra is the closest relative to Motivo's compositional ambitions, since both treat repetition and
variation as first-class, composable operations rather than as duplicated note data.
Strudel ports that model to the browser with WebAudio, OSC, and optional MIDI~\cite{Roos_McLean_2023}, broadening reach
without changing the underlying execution model.
Both systems evaluate patterns at show time and stream events to external backends; even MIDI output serves a live
session rather than a deterministic compile pass.
Motivo instead targets a self-contained Format~1 file importable without the pattern engine running, so the timeline
exists independently of any performance clock.

\topic{HMusic.} Compiles Haskell-embedded track and pattern descriptions to Sonic Pi~\cite{Du_Bois_Ribeiro_2019},
representing music as grid-like patterns and parallel tracks that can be combined with algebraic operators.
It shares Motivo's interest in a compilation step and in abstracting repetitive structure above the level of individual
events.
Its target, however, is Sonic Pi source, so the generated program inherits the same live-synthesis assumptions: the
output remains performance-coupled rather than a portable symbolic artifact validated independently of a running
backend.

\topic{ChucK.} Foregrounds sample-level synthesis with a \emph{strongly-timed} model in which the program advances
time explicitly and synchronizes concurrent ``shreds''~\cite{Wang_Cook_2003,Wang_2008}.
This gives composers fine-grained, deterministic control over \emph{when} computation happens, an interest Motivo
shares in spirit through its compile-time scheduling of events.
The control, however, operates at the level of audio time and concurrent processes rather than hierarchical song
structure; high-level compositional constructs such as motifs, sections, and reusable patterns are not the unit of
abstraction, and there is no offline pass that emits a portable score.

\topic{SuperCollider.} Separates a synthesis server from a client language and expresses music through unit generators,
patterns, and streams driven over OSC~\cite{McCartney_2002,Wilson_Cottle_Collins_2011}.
Its pattern and stream libraries demonstrate, like Motivo, that compact descriptions can generate large amounts of
musical material.
Yet composition and execution are not separated: patterns are realized by a running server, so the description and its
playback are bound together, and there is no notion of compiling a complete piece down to a standalone symbolic file.
Across all of these systems, composition and execution tend to be intertwined.
Motivo instead compiles source to a self-contained Format~1 MIDI file with no Motivo runtime at playback time.
